\documentclass[11pt]{beamer}
%\usetheme{Frankfurt}
%\usetheme{Singapore}
\usetheme{CambridgeUS}
\usecolortheme{whale}
\definecolor{azulescuro}{RGB}{0, 32, 96}
\setbeamercolor{frametitle}{fg=azulescuro}
\setbeamertemplate{navigation symbols}{}
\usepackage[utf8]{inputenc}
\usepackage{amsmath}
\usepackage{amsfonts}
\usepackage{amssymb}
\usepackage{graphicx}
\usepackage{multicol}
\usepackage{changepage}
\usepackage[labelformat=empty]{caption}
\usepackage{booktabs}
%\usepackage{natbib}
\usepackage{bibentry}
\usepackage{url}
%\usepackage{enumitem}
%\bibliographystyle{alpha}
\nobibliography *
\usepackage[brazil]{babel}
%\usepackage{abnt-alf}
%\usepackage[alf]{abntcite}
\author[UFCG]{Beatriz Sabino, Bianka Abreu, M. Eduarda Sinfrônio}
\title{Energia renovável e desenvolvimento econômico no Nordeste}
%\setbeamercovered{transparent} 
%\setbeamertemplate{navigation symbols}{} 
\institute{UFCG-UAEF} 
\date{\today}
\makeindex
 
\begin{document}
\begin{frame}
\titlepage
\end{frame}
\begin{frame}{Sumário}
%\begin{multicols}{2}
		\centering
		\tableofcontents
%\end{multicols}
\end{frame}
\begin{frame}{Introdução}
\section{Introdução}
\begin{itemize}
    \item Mudanças climáticas e altas emissões de GEE impulsionam a busca global por energias limpas
    \item Estudos (Grécia, EUA) mostram que renováveis impulsionam PIB e geração de empregos
    \item Brasil possui uma das matrizes elétricas mais limpas do mundo, historicamente hidrelétrica
    \item Energia eólica: 1ª turbina em Fernando de Noronha (1992); crise de 2001 e PROINFA (2002) consolidam o setor
    \item Hoje, a eólica responde por \textbf{14\%} da matriz nacional -- 2ª maior fonte elétrica do país
\end{itemize}
\end{frame}
\begin{frame}{Introdução}
\begin{itemize}
   \item Nordeste concentra \textbf{91\%} dos parques eólicos do Brasil, devido aos ventos alísios constantes e em picos de vento, a eólica chega a suprir \textbf{60\%} da demanda elétrica da região 
    \item BNB (FNE Sol) já financiou bilhões de reais em projetos de energia limpa na região
    \item \textbf{Objetivo}: investigar a relação entre expansão eólica e desenvolvimento socioeconômico nos 9 estados do NE
    \item Painel de dados (2010--2024): potência acumulada em MW (ANEEL) x IDHM
\end{itemize}
\end{frame}
\begin{frame}
\frametitle{Referencial Teórico}
\section{Referencial Teórico}
\begin{itemize}
    \item Desenvolvimento econômico
    \item Matriz elétrica brasileira        
    \item Energia renovável 
    \item Nordeste

\end{itemize}
\end{frame}
\begin{frame}
\frametitle{Metodologia}
\section{Metodologia}
\begin{itemize}
    \item Natureza da pesquisa
    \item Base de dados
    \item Variáveis
    \item Modelo estatístico
    \item Software R
\end{itemize}
\end{frame}
\begin{frame}
\frametitle{Análise de Dados e Discussão dos Resultados}
\section{Análise de Dados e Discussão dos Resultados}
\framesubtitle{Estatística Descritiva}
\begin{itemize}
    \item Painel com \textbf{126 observações} (9 estados, 14 anos: 2010--2024)
    \item IDHM: média de \textbf{0,71}, com baixíssima variação entre estados (desvio padrão de 0,03) -- indica convergência regional
    \item Potência Acumulada (MW): média de \textbf{2.299 MW}, mas desvio padrão de 2.562 MW -- maior que a própria média, revelando forte desigualdade
    \item 25\% das observações têm capacidade quase nula, enquanto o topo chega a 13.620 MW (Bahia e Rio Grande do Norte)
    \end{itemize}
\end{frame}
\begin{frame}
\frametitle{Análise de Dados e Discussão dos Resultados}
\framesubtitle{Evolução Regional (2010 vs. 2024)}
\begin{itemize}
     \item Em 2010, infraestrutura quase inexistente: Piauí e Alagoas com menos de 1 MW; Sergipe sem geração registrada
    \item Em 2024, Bahia lidera com \textbf{13.620 MW} (era 325 MW), seguida do Rio Grande do Norte com \textbf{11.288 MW}
    \item Piauí tem o maior crescimento relativo, alcançando 6.490 MW
    \item Todos os 9 estados saem da faixa de IDHM ``médio'' (2010) para ``alto'' (2024), com destaque para RN e CE
    \end{itemize}
\end{frame}
\begin{frame}
\frametitle{Análise de Dados e Discussão dos Resultados}
\framesubtitle{Boxplots e Tendências}
\begin{itemize}
    \item Crescimento da capacidade acelera a partir de \textbf{2018}, puxado por Bahia e Rio Grande do Norte
    \item Boxplots confirmam a desigualdade: BA e RN com caixas longas e alta variabilidade
    \item Alagoas mantém investimentos baixos e estáveis ao longo da série
    \item IDHM tem distribuição mais compacta e uniforme -- avanço social mais homogêneo que a expansão energética
    \end{itemize}
\end{frame}
\begin{frame}
\frametitle{Análise de Dados e Discussão dos Resultados}
\framesubtitle{Correlação entre Energia e Desenvolvimento}
\begin{itemize}
    \item Gráfico de dispersão mostra relação positiva clara entre potência acumulada e IDHM
    \item Correlação de Pearson: $r = 0,59$ -- correlação moderada a forte
    \item Expansão da energia renovável no Nordeste caminha lado a lado com a melhoria dos indicadores sociais (2010--2024)
    \end{itemize}
\end{frame}
\begin{frame}
\frametitle{Considerações Finais}
\section{Considerações Finais}
\begin{itemize}
    \item \textbf{Objetivo}: investigar a relação entre expansão da energia renovável e desenvolvimento socioeconômico no Nordeste (painel de 9 estados, 2010--2024: potência em MW x IDHM)
    \item Expansão energética \textbf{irregular}, concentrada em Bahia e Rio Grande do Norte; IDHM cresceu de forma \textbf{generalizada}, com todos os estados migrando de desenvolvimento médio para alto
    \item Correlação de Pearson de \textbf{0,59} entre potência instalada e IDHM: associação moderada, mas que não implica causalidade -- políticas públicas e transferências governamentais também influenciam
\end{itemize}
\end{frame}
\begin{frame}
\frametitle{Considerações Finais}
\begin{itemize}
    \item \textbf{Limitações}: apenas duas variáveis analisadas, sem testes de causa e efeito; recomenda-se, em pesquisas futuras, dados em painel com variáveis de controle e novos indicadores (emprego, renda, arrecadação)
    \item Energia renovável é estratégica para o Nordeste, mas seu potencial de desenvolvimento depende de políticas que reduzam desigualdades na distribuição dos investimentos entre os estados
    \end{itemize}
\end{frame}
\begin{frame}[allowframebreaks,c]
\frametitle{Referências}
\section{Referências}
\footnotesize
\begin{itemize}
    \item BARBOSA, Wladimir De Macedo. \textit{Impactos econômicos das energias renováveis na economia nordestina}. Repositório institucional, 2024. Disponível em: \url{https://repositorio.ufc.br/bitstream/riufc/76832/1/2024_dis_wmbarbosa.pdf}. Acesso em: 01 jun. 2026.
    \item DE OLIVEIRA, Magnus Kelly; UHR, Daniel De Abreu Pereira; UHR, Júlia Gallego Ziero. \textit{Ventos de mudança: avaliando o impacto econômico dos parques eólicos nos municípios do Nordeste brasileiro}. Revista Econômica do Nordeste, 2025. Disponível em: \url{https://www.bnb.gov.br/revista/ren/article/view/2010/2320}. Acesso em: 10 jun. 2026.
    \item MACRO, Do Micro Ao. \textit{Banco do Nordeste anuncia R\$ 360 milhões para ampliar rede de energia renovável no país}. Carta Capital, 2025. Disponível em: \url{https://www.cartacapital.com.br/do-micro-ao-macro/banco-do-nordeste-anuncia-r-360-milhoes-para-ampliar-rede-de-energia-renovavel-no-pais/}. Acesso em: 10 jun. 2026.
    \item ENERGIA, Neo. \textit{Energia eólica: ventos do Nordeste}. Neoenergia, 2020. Disponível em: \url{https://www.neoenergia.com/w/energia-eolica-ventos-do-nordeste}. Acesso em: 11 jun. 2026.
    \item FREIRE, Anderson Ítalo; FONTGALLAND, Isabel Lausanne. \textit{Perspectives and economic challenges of wind energy generation in the Northeastern region of Brazil}. Research, Society and Development, 2022. Disponível em: \url{https://rsdjournal.org/rsd/article/view/25429/22221}. Acesso em: 11 jun. 2026.
\end{itemize}
\bibliographystyle{alpha}
\bibliography{bib/ipd_game.bib}
\end{frame}
\end{document}